% ==============================================================================
% CONCLUSION GÉNÉRALE — NG-IDPS
% Système Autonome de Détection et Prévention d'Intrusions Nouvelle Génération
% Master Professionnel : Computer Engineering — Expert Réseaux, ISIMa
% ==============================================================================
% Comme pour l'introduction, le titre de chapitre est posé dans le fichier principal :
% \chapter*{Conclusion Générale}
% \addcontentsline{toc}{}{\hspace*{-1.5em}\textbf{Conclusion Générale}}
\label{chap:conclusion}

Ce projet de fin d'études avait pour ambition de concevoir, de réaliser et de valider un système de détection et de prévention des intrusions réseau capable de fonctionner de manière autonome dans un environnement local et isolé, au profit d'organisations qui ne disposent ni d'équipes de supervision permanentes ni de plateformes commerciales intégrées. Le système NG-IDPS qui en résulte associe, au sein d'une chaîne unique, une détection par signatures, une détection comportementale non supervisée, un renseignement sur les menaces, une orchestration de la réponse et un tableau de bord de supervision, le tout reposant sur des composants libres et sur des modèles de langage auto-hébergés pour ses fonctions critiques.

La phase préparatoire a permis de traduire la problématique en sept objectifs fonctionnels et cinq objectifs non fonctionnels, de planifier le travail en six sprints selon la méthodologie Scrum et de définir une architecture répartie en nœuds spécialisés : capture, apprentissage, réponse et supervision.

La première \textit{Release}, consacrée à la détection, a mis en place la capture du trafic par une sonde Suricata, un pipeline à double chemin (indexation pour l'investigation, flux temps réel pour l'analyse) et un moteur de détection comportementale fondé sur un autoencodeur variationnel à mémoire longue (LSTM-VAE). Ce dernier fonctionne avec une latence d'inférence de 2 à 4 millisecondes par séquence, très en deçà de la cible de 15 millisecondes. Son seuil a été calibré en comparant quatre points de fonctionnement ; le point retenu produit un rappel de 51,89\,\% pour 9,23\,\% de séquences signalées à tort, une aire sous la courbe ROC de 0,79 et une aire sous la courbe précision-rappel de 0,84. Ces faux positifs sont absorbés par une règle de persistance temporelle avant toute action automatique.

La seconde \textit{Release}, consacrée à la prévention et à la réponse, a ajouté quatre capacités : une veille sur les menaces (CTI) qui agrège des sources ouvertes, trie les articles et propose des règles de détection candidates soumises à validation ; une chaîne de sensibilisation qui transforme les menaces à caractère humain en parcours de formation ciblés, avec un suivi de conformité ; un orchestrateur de réponse à plusieurs paliers, capable d'isoler une source malveillante de façon autonome ou après validation d'un analyste ; et un tableau de bord SOC présentant les incidents, leur diagnostic structuré par un modèle de langage local, l'état de l'infrastructure et un rapport d'audit exportable.

La validation ne s'est pas limitée à des mesures sur un jeu de données de référence. Le système a été confronté à des attaques réellement exécutées depuis une machine dédiée (balayages de ports, identification de services, force brute SSH, scan applicatif web, inondation SYN), dont les six campagnes rejouées dans l'ordre chronologique ont toutes été détectées, avec un taux de $0{,}18\,\%$ de fenêtres bénignes signalées à tort sur le trafic réel des nœuds de l'infrastructure. Cette démarche a surtout permis d'identifier puis de corriger des défauts concrets que l'implémentation seule ne révélait pas : une sélection de caractéristiques défavorable à certaines classes d'attaques, un fichier de règles non référencé par la sonde, des faux positifs provoqués par le trafic du pipeline lui-même, et, du côté des agents d'intelligence artificielle, des sorties désorganisées, des règles recopiées d'un jeu de référence, des étiquettes de source qui biaisaient le classement et des droits d'accès insuffisants pour l'agent de déploiement.

Ces résultats permettent de répondre aux cinq défis posés en introduction. Le défi de la latence est atteint : l'inférence compilée, le tampon glissant de dix pas et les pré-filtres volumétriques maintiennent le temps de calcul entre 2 et 4 millisecondes par séquence, pour une décision d'isolation émise environ dix secondes après le début des attaques testées, délai dominé par le remplissage de la fenêtre d'observation plutôt que par le calcul lui-même. Le défi des faux positifs est également atteint : la calibration du seuil, la règle de persistance sur cinq évaluations et la neutralisation des actifs connus par liste blanche ramènent le taux de séquences signalées à tort à 9,23\,\% hors ligne et à 0,18\,\% sur le trafic réel de l'infrastructure. Le défi de la détection de comportements inconnus est atteint grâce à la complémentarité entre la détection par signatures et la détection non supervisée, qui couvrent ensemble aussi bien les attaques connues que les comportements inédits. Le défi de l'accessibilité est atteint : l'emploi de composants libres et d'un modèle de langage local pour l'analyse des incidents et la veille, les modèles distants n'étant sollicités que pour la génération de contenus pédagogiques et toujours avec un repli local, rend la chaîne complète exécutable sans licence ni service obligatoire. Le défi de la continuité de la réponse, enfin, est atteint : l'orchestrateur à plusieurs paliers isole une source malveillante de façon autonome en l'absence d'opérateur, produit des rapports d'audit exportables et détecte la panne d'un nœud en moins de quinze secondes grâce à la surveillance continue de l'infrastructure.

Au regard des objectifs fixés au premier chapitre, les sept objectifs fonctionnels ont tous été réalisés. Parmi les objectifs non fonctionnels, la latence d'inférence, l'autonomie opérationnelle et la maîtrise de l'infrastructure sont pleinement atteintes. La performance de détection et la fiabilité du système le sont dans les limites de l'évaluation menée, que les paragraphes suivants précisent.

Ce travail comporte plusieurs limites qu'il convient d'assumer. La première concerne la portée de l'évaluation en conditions réelles : les attaques réelles ne sont qu'au nombre de six, exécutées dans un laboratoire à site unique, et un taux de détection de 6 sur 6 reste statistiquement compatible, à 95\,\% de confiance, avec un taux réel compris entre 61 et 100\,\%. Le jeu de données de référence, issu de 2017, ne reflète par ailleurs pas toutes les pratiques d'attaque actuelles ; ces résultats attestent donc le fonctionnement de bout en bout de la chaîne, sans valeur statistique généralisable. La deuxième limite tient à un écart entre entraînement et production : certaines caractéristiques, notamment les drapeaux TCP, sont approximées en production à partir des agrégats de flux et non comptées paquet par paquet comme dans les données d'entraînement. Le dispositif compense cet écart par une neutralisation réservée aux adresses de la liste blanche, de sorte que le comportement du modèle sur un hôte inconnu au trafic ordinaire n'a pas été évalué, et que le rappel hors ligne reste limité sur le balayage de ports. La troisième limite porte sur la fiabilité des modèles de langage de petite taille : le modèle local de trois milliards de paramètres produit des diagnostics utiles mais perfectibles, avec parfois une technique MITRE ATT\&CK inexacte, des recommandations de remédiation inégales, des scores de fiabilité aberrants ou des règles recopiées d'un jeu de référence ; ces sorties demeurent des propositions, et la validation humaine, le contrôle de structure ainsi que le test à blanc par le moteur Suricata restent indispensables, ce dernier ne vérifiant au demeurant que la syntaxe d'une règle et non sa pertinence. La quatrième limite concerne la sécurité du plan de contrôle interne : dans le laboratoire isolé où le système a été développé, les commandes de remédiation transmises par le canal temps réel ne sont pas authentifiées et le bus de messages n'est ni authentifié ni chiffré, une configuration qui n'est pas acceptable hors d'un réseau fermé. La dernière limite touche au périmètre réseau et à la capacité du système : le blocage repose sur l'adresse IPv4 source, les adresses IPv6 de portée locale étant écartées, la protection de la couche liaison ne couvre qu'un nœud et qu'une adresse, et le passage à l'échelle, qu'il s'agisse d'un courtier de messages unique ou d'une inférence exécutée sur processeur, n'a pas été évalué.

Plusieurs évolutions permettraient de prolonger ce travail. La première consisterait à rétablir la parité entre entraînement et production, en calculant de véritables compteurs de drapeaux à la périphérie ou en recalibrant et en réentraînant le modèle sur du trafic bénin local capturé avec le même extracteur, ce qui rendrait la neutralisation actuelle inutile. La deuxième consisterait à élargir l'évaluation à des campagnes plus nombreuses et plus variées, incluant des attaques lentes, des tentatives d'évasion et un trafic bénin diversifié, afin d'obtenir des intervalles de confiance resserrés. La troisième consisterait à renforcer la validation des règles générées par un test sur trafic rejoué avant déploiement, en complément du test à blanc syntaxique déjà en place, ainsi que par un contrôle de leur pertinence vis-à-vis de l'article source. La quatrième consisterait à sécuriser le plan de contrôle par l'authentification des clients et des commandes, le chiffrement et l'authentification du bus de messages, et la signature des règles déployées. La dernière consisterait à exploiter les décisions des analystes, qu'il s'agisse d'approbations ou de faux positifs déclarés, pour recalibrer périodiquement le seuil de détection et réentraîner le modèle.

Sur le plan personnel, ce projet nous a conduit à articuler des compétences habituellement séparées : réseaux et sécurité, apprentissage automatique, ingénierie logicielle distribuée et conduite de projet agile. Il nous a surtout appris que la valeur d'un système de sécurité se mesure autant à la rigueur avec laquelle on le met à l'épreuve qu'à ses performances annoncées. En ce sens, NG-IDPS constitue une base crédible et améliorable pour les organisations qui cherchent à acquérir une capacité de cyberdéfense proactive tout en conservant la maîtrise de leurs données.